\documentclass[aps,prl,twocolumn,superscriptaddress]{revtex4-2}
\usepackage{amsmath,amssymb,graphicx,hyperref}

\begin{document}

\title{BCT Appendix App-L207b:\\
The Great Pyramid Slope Angle as a BCT Void Projection}
\author{Michel Robert Cabri\'e}
\email{ZeroFreeParameters@gmail.com}
\affiliation{Independent Artist and Researcher, Victoria, Australia\\
ORCID: 0009-0007-9561-9859}
\date{March 2026}

\begin{abstract}
We derive the slope angle of the Great Pyramid of Giza from BCT void geometry. The observed angle of $51.84^\circ \pm 0.05^\circ$ matches $\arctan(4/\pi) = 51.854^\circ$ to within $0.03\%$. The ratio $4/\pi$ arises naturally as the projection of the BCT octahedral void onto a two-dimensional surface: the ratio of the void's circumference ($2\pi r_{\mathrm{oct}}$) to its projected diameter ($4 r_{\mathrm{oct}} \times \pi / 4$) produces the observed slope. This is the same $\pi$-encoding that appears in the pyramid's base-to-height ratio: $2 \times \text{base} / \text{height} = \pi$ to within $0.04\%$. Zero free parameters.
\end{abstract}

\maketitle

\section{The Observed Angle}

The Great Pyramid (Khufu) has a measured slope angle of $51^\circ 50' 40'' \pm 1'$, equivalent to $51.844^\circ \pm 0.017^\circ$~\cite{Petrie1883, Lehner1997}.

The base length is $230.33 \pm 0.05$~m and the original height was $146.59 \pm 0.20$~m, giving:
\begin{equation}
\theta_{\mathrm{obs}} = \arctan\left(\frac{2 \times 146.59}{230.33}\right) = \arctan(1.2728) = 51.844^\circ
\end{equation}

\section{The BCT Derivation}

The ratio $4/\pi$ appears in BCT when the octahedral void is projected from three dimensions onto a two-dimensional surface:

The octahedral void in the BCT lattice has radius $r_{\mathrm{oct}} = (\sqrt{2}-1)/2$. Its three-dimensional surface area is:
\begin{equation}
A_{3D} = 4\pi r_{\mathrm{oct}}^2
\end{equation}

Its two-dimensional projected area (shadow on a plane) is:
\begin{equation}
A_{2D} = \pi r_{\mathrm{oct}}^2
\end{equation}

The ratio:
\begin{equation}
\frac{A_{3D}}{A_{2D}} = 4
\end{equation}

When this dimensionless ratio is used as the tangent argument with the $\pi$-normalisation natural to circular projection:
\begin{equation}
\theta_{\mathrm{BCT}} = \arctan\left(\frac{4}{\pi}\right) = 51.854^\circ
\end{equation}

\textbf{Error:} $51.854 - 51.844 = 0.010^\circ$, or $+0.019\%$.

\section{The $\pi$-Encoding}

The same geometry produces the well-known $\pi$-relationship in the pyramid's dimensions:
\begin{equation}
\frac{2 \times \text{base}}{\text{height}} = \frac{2 \times 230.33}{146.59} = 3.1427
\end{equation}

Compare: $\pi = 3.14159$. Error: $+0.04\%$.

In BCT, this is not a coincidence but a geometric necessity: a structure built with slope angle $\arctan(4/\pi)$ automatically encodes $\pi$ in its base-to-height ratio. The builders did not need to ``know'' $\pi$ explicitly --- they needed only to replicate the angle of the octahedral void projection.

\section{How Would They Know?}

This is the question that ancient mystery enthusiasts and professional Egyptologists both ask. BCT offers a specific answer: the angle is the projection of the most natural geometric void in a sphere packing.

If the builders were working with physical sphere models --- stacking balls, observing the gaps --- the octahedral void projection angle would emerge naturally from the geometry of the stack. No advanced mathematics required. No alien intervention. Just the geometry of touching spheres, which produces $\arctan(4/\pi)$ as inevitably as it produces $r_{\mathrm{oct}} = (\sqrt{2}-1)/2$.

The builders were, in this reading, the first BCT engineers --- reading vacuum geometry from sphere packing models, millennia before the lattice had a name.

\section{Other Pyramids}

The Red Pyramid at Dahshur has slope angle $43.36^\circ$. The Bent Pyramid transitions from $54.27^\circ$ to $43.22^\circ$. These angles correspond to different BCT void projections:
\begin{align}
\arctan\left(\frac{r_{\mathrm{oct}}}{r_{\mathrm{tet}}}\right) &= \arctan(1.843) = 61.53^\circ \\
\arctan\left(\frac{2}{\pi}\right) &= 32.48^\circ \\
\arctan(1) &= 45.00^\circ
\end{align}

The observed angles do not match these simple BCT ratios as cleanly as the Great Pyramid matches $\arctan(4/\pi)$. The Great Pyramid's angle is the most precisely BCT-aligned of all Egyptian pyramids --- consistent with it being the culmination of a geometric tradition.

\begin{acknowledgments}
Sir Flinders Petrie measured the Great Pyramid to sub-arcminute precision in 1881--1882. His measurements, made with theodolites and steel tapes in the Egyptian heat, remain among the most accurate ever taken. This appendix is built on his data.
\end{acknowledgments}

\begin{thebibliography}{4}
\bibitem{Petrie1883} W.~M.~F.~Petrie, The Pyramids and Temples of Gizeh (Field \& Tuer, 1883).
\bibitem{Lehner1997} M.~Lehner, The Complete Pyramids (Thames \& Hudson, 1997).
\bibitem{L207} M.~R.~Cabri\'e, BCT Letter 207: The Ancient Geometers, Zenodo (2026).
\bibitem{L189} M.~R.~Cabri\'e, BCT Letter 189: BCT-SAR Archaeology, Zenodo (2026).
\end{thebibliography}

\end{document}
